\bookchapter{第零章：嗨，開始之前}

在正式探索前，咱們先聊點基本規則，確保後面說的你都懂。別急，磨刀不誤砍柴工嘛。

\booksection{一、這本書能幫你做什麼？（和不能做的）}

它能幫你系統性地搞懂手機怎麼用，明白功能背後的道理，而且無論你是安卓黨還是蘋果派，大多數思路是通用的。它會帶著你動手練習，讓你從“會用”升級到“懂用”。

它不是某一款特定手機的說明書，也不是追新品的科技新聞，更不是教你沉迷手機的指南。說實話，如果你期待的是“如何用手機刷影片更爽”，那可能要失望了——我們更關心的是如何讓手機真正為你服務，而不是反過來。

\booksection{二、我們的“暗號”說明}

為了說得更清楚，書裡會用一些小標記，這些標記就像我們之間的默契暗號：

\begin{itemize}
  \item \textbf{[按鈕名]}：代表螢幕上具體的按鈕或App名字，比如[設定]、[確定]。
  \item \textbf{>}：代表操作順序，比如"點 [設定] > [聲音]"，就是先點設定，再點聲音。
  \item \textbf{加粗}：非常重要！可能是關鍵概念或安全警告，請務必留意。
  \item \emph{斜體}：一些補充的小想法或比喻，幫你更好地理解。
\end{itemize}

\booksection{三、安卓和iOS，咱都說說}

你猜怎麼著，有些東西叫法不同，但乾的事兒一樣。比如：應用（安卓）和 App（iOS），設定（都叫設定），檔案管理（安卓）和 檔案（iOS）。知道是一回事就行，別被名詞繞暈了。

另外，現在安卓陣營裡還有華為的鴻蒙系統，它有自己的特色，後面我們會專門講到。總之，不管你用什麼手機，這本書都能幫到你。

\booksection{四、手機容量：為啥買了128G，可用卻不到110G？}

容量單位大家知道：1TB = 1024GB，1GB = 1024MB，1MB = 1024KB……但這裡有個坑：硬碟廠商通常按1000進位制算，而系統按1024進位制算，所以你買到的“128G”硬碟，實際容量會“縮水”。

但你手機容量縮水的原因主要有三個：

\begin{enumerate}
  \item \textbf{系統自己佔了地方}：作業系統就像房子的承重牆，必須存在。現在的系統動輒佔用10-20GB，這很正常。
  \item \textbf{廠商預裝的應用}：新手機裡自帶的那些App也佔空間，有的能解除安裝，有的不能。
  \item \textbf{分割區和格式化損耗}：技術原因，就像書包的襯布也佔體積。
\end{enumerate}

\begin{tcbinfo}{買手機的小建議}
現在128GB的手機基本不夠用了，256GB起步比較穩妥。如果你喜歡拍照、玩遊戲，建議直接上512GB甚至更大。記住，儲存空間這東西，寧可富餘，不可捉襟見肘。

\end{tcbinfo}

\begin{tcbpractice}{來動手：看看你的手機"身份證"}
開啟 [設定]，找到關於手機（或類似選項），記下：手機型號、系統版本、總儲存空間和剩餘空間。這是瞭解它的第一步。順便看看電池健康度，心裡有個數。

\end{tcbpractice}

\booksection{五、你的手機是一座“數字資產庫”}

在開始具體知識之前，我們先建立一個核心認知：\textbf{手機早已不只是通訊工具，它是我們最重要的數字資產庫}。聯絡人、聊天記錄、家庭照片、工作文件、銀行App、付款碼、健康資料、電子證件……這些資料加在一起，價值往往超過手機本身。

\begin{tcbinfo}{為什麼這很重要？}
如果手機丟了或壞了，你損失的不僅是幾千塊錢的硬體，\textbf{更可能是多年積累的不可再生資料}——孩子的成長照片、過世親人的語音、沒備份的重要合同。科學地整理與防護習慣，能從三個維度保護你：

\begin{itemize}
  \item \textbf{提升執行效率}：規範的檔案分類讓系統查詢更快，手機長期不卡頓。
  \item \textbf{延長使用壽命}：合理的儲存管理減少晶片讀寫負擔，電池和儲存都更耐用。
  \item \textbf{保障資料安全}：分層管理 + 定期備份 + 許可權收斂，讓"資料災難"遠離你。
\end{itemize}

\end{tcbinfo}

\booksubsection{我們的核心理念：“預防為主，防治結合”}

就像養生不能等到得病才看醫生，\textbf{手機管理也要從“日常習慣”抓起}。本書反覆強調的不是某個高階技巧，而是這些樸素的原則：

\begin{itemize}
  \item \textbf{預防為主}：裝 App 只走官方商店、照片發之前先脫敏、密碼不重複——這些"小動作"成本極低，但能擋住絕大多數風險。
  \item \textbf{定期體檢}：每月一次整理下載資料夾、每季度一次許可權審計、每半年一次全量備份——固定節奏比臨時救火有效一百倍。
  \item \textbf{留有退路}：開啟雲同步、記好賬號密碼、寫一份"數字遺產清單"——確保意外發生時你和家人都能從容應對。
\end{itemize}
